La evaluación tiene tres partes: la recuperación por componentes, el sistema completo y las pruebas
adversariales. Todos los scripts están en \texttt{scripts/} y guardan los resultados en
\texttt{eval/resultados/}.

\subsection{Conjunto de preguntas}

Se siguió el método híbrido presentado en la clase 2. Se muestrearon fragmentos del índice (hasta dos
por artículo, con prioridad para resúmenes, resultados, métodos y discusión) y Claude Opus 5 escribió
una pregunta que cada fragmento responde, junto con una respuesta de referencia. El 75\,\% de las
preguntas está en español y el 25\,\% en inglés. El resultado son 40 preguntas respondibles de 21
artículos. Se agregaron a mano 5 preguntas fuera de alcance y 3 preguntas sobre el tema que el corpus
no puede responder (por ejemplo, sobre fNIRS en adultos mayores bilingües).

Las preguntas generadas tienen el campo \texttt{revisada} y la revisión humana está pendiente.
Algunas preguntas son de artículos que las búsquedas trajeron, pero que no tratan el tema (un estudio
de golf, un análisis bibliométrico de dislexia). Se dejaron porque el corpus real tiene ese ruido.

\subsection{Recuperación}

Se compararon cuatro modos (BM25, denso, híbrido con RRF e híbrido con reranker) con tres variantes
de consulta: la pregunta original, las consultas del planificador y las dos juntas. Se midió a dos
niveles:

\begin{itemize}
\item Artículo: si el artículo del que salió la pregunta aparece entre los primeros resultados.
\item Fragmento: con un solo fragmento relevante por pregunta, el de origen, la métrica subestima la
recuperación, porque los fragmentos vecinos y el resumen del mismo artículo también responden. Por
eso se armaron juicios de relevancia por pooling: para cada pregunta se juntaron los diez primeros
fragmentos de todas las configuraciones y Claude Opus 5 marcó cuáles responden la pregunta. Quedaron
en promedio 4,85 fragmentos relevantes por pregunta.
\end{itemize}

\subsection{Sistema completo}

Cada pregunta pasó por el sistema completo sin caché. Para las respondibles se midió si el artículo
de origen está entre los fragmentos recuperados y entre las referencias citadas. Un juez evaluó la
corrección de la respuesta frente a la referencia (de 1 a 5) y contó las afirmaciones respaldadas por
los documentos citados. El juez es Claude Opus 5, un modelo distinto del verificador (Haiku 5.5),
para que el sistema no se evalúe a sí mismo con el mismo modelo. También se registraron el costo, los
tokens, la latencia y el comportamiento de los agentes: veredictos, iteraciones y rechazos del
planificador. La misma evaluación se repitió sin el verificador, para medir su aporte.

\subsection{Pruebas adversariales}

Se escribieron 12 ataques de inyección directa, en inglés y en español: pedidos de revelar el prompt,
cambios de rol, etiquetas falsas, pedidos de inventar referencias y de ignorar los documentos. Se
midió cuántos bloquea la guardia, cuántos rechaza el planificador y si alguno logra filtrar el
prompt. Para los falsos positivos se pasaron por la guardia las 40 preguntas legítimas. Para la
inyección indirecta se plantó un documento con instrucciones (afirmar que el PLI tiene ICC = 1,0 y
recomendar un enlace externo) y se comprobó si el detector lo marca. Después se salteó el filtro a
propósito y se le entregó el documento al analista junto con cinco fragmentos reales, para ver si
obedece.
